% Generated draft: requires mathematical and source review.
\documentclass[11pt]{article}
\usepackage[margin=25mm]{geometry}
\usepackage{fontspec}
\setmainfont{Latin Modern Roman}
\usepackage{amsmath,amssymb,mathtools,unicode-math}
\setmathfont{Latin Modern Math}
\usepackage{enumitem,needspace,xurl,hyperref}
\setlength{\parindent}{0pt}
\setlength{\parskip}{5pt}
\setlength{\emergencystretch}{3em}
\newcommand{\N}{\mathbb{N}}
\newcommand{\Z}{\mathbb{Z}}
\newcommand{\Q}{\mathbb{Q}}
\newcommand{\R}{\mathbb{R}}
\newcommand{\C}{\mathbb{C}}
\newcommand{\F}{\mathbb{F}}
\newcommand{\A}{\mathbb{A}}
\newcommand{\PP}{\mathbb{P}}
\newcommand{\E}{\mathbb{E}}
\newcommand{\eps}{\varepsilon}
\newcommand{\dd}{\,\mathrm{d}}
\begin{document}
\section*{UnsolvedMath ID: 700655}
\textbf{Source problem code:} catalogue entry
\subsection{Definitions and mathematical statement}
Let \(G\) be a group with identity \(1_G\). Its integral group ring \(\Z G\) consists of finite sums \(\sum_{g\in G}a_g g\), where \(a_g\in\Z\). For a positive integer \(n\), let \(e\in M_n(\Z G)\) be idempotent, meaning \(e^2=e\), and write its ordinary matrix trace as
\[
\operatorname{tr}(e)=\sum_{i=1}^{n}e_{ii}=\sum_{g\in G}e_g g,\qquad e_g\in\Z.
\]
Only finitely many coefficients \(e_g\) are nonzero. Write \(g\sim x\) if \(g=hxh^{-1}\) for some \(h\in G\).

The supplied statement attributes the following conjectures to H. Bass, with universal quantification over all groups \(G\), all positive integers \(n\), and all idempotents \(e\in M_n(\Z G)\):
\begin{enumerate}
\item \textbf{Strong conjecture.} For every \(x\in G\setminus\{1_G\}\),
\[
\sum_{\substack{g\in G\\g\sim x}}e_g=0.
\]
\item \textbf{Weak conjecture, as transcribed in the supplied record.}
\[
\sum_{g\in G}e_g=0.
\]
\end{enumerate}
The weak formulation has an unresolved source defect: the sum explicitly includes \(1_G\), and no normalization or restriction excluding the identity idempotent is supplied. For \(e=(1_{\Z G})\), that sum equals \(1\). Thus the literal weak formulation is incompatible with the accompanying assertion that it holds for residually finite groups. The intended weak conjecture cannot be recovered unambiguously from this record; no correction is imposed here.

The supplied source reports that the strong conjecture has been proved for finite, abelian, and linear groups, and that the weak conjecture has been proved for residually finite groups. These are source-reported qualifications, not independently verified assessments; the latter assertion is subject to the unresolved weak-formulation defect above.
\subsection{Short English statement}
For every idempotent matrix over an integral group ring, the strong Bass conjecture asks whether the trace coefficients sum to zero on each nonidentity conjugacy class. The supplied weak formulation asks for the sum over the entire group to vanish, but this conflicts with the allowed identity idempotent and requires source clarification. The registry labels the problem open; this label has not been independently verified.
\subsection{Sources}
\begin{itemize}
\item Local UnsolvedMath registry, supplied complete source record, permanent ID \(700655\), exact source problem code \texttt{catalogue entry}, title \emph{Kourovka Notebook Problem 10.40}. The supplied registry status is \texttt{open}; it is not an independently verified mathematical assessment. Upstream revision: \texttt{d371ee917ac9647fbb06cd6fc529deec6f3dda82}.
\item \emph{The Kourovka Notebook (Unsolved Problems in Group Theory)}, original problem 10.40, 10th issue (1986); supplied edition reference: 21st issue edition (2026), \texttt{arXiv:1401.0300v48}. The supplied collection description names E. I. Khukhro and V. D. Mazurov as editors. The record names G. A. Noskov as proposer and attributes the conjectures in the statement to H. Bass.
\item H. Bass's comment of 2005, as reproduced in the supplied record, reports significant progress and recommends the following paper as an account current at that time: A. J. Berrick, I. Chatterji, and G. Mislin, \emph{Math. Ann.} \textbf{329} (2004), 597--621. The supplied record gives no paper title or link. This historical assessment and the bibliographic reference have not been independently verified.
\end{itemize}
\end{document}
